\documentclass[../../../../main.tex]{subfiles}
\begin{document}

\begin{flushright}
\footnotesize\textit{【自動文字起こし・要確認】}
\end{flushright}

\noindent
\textbf{[解]} $NA, NB$ と $k$ の交点を $C, D$ とする。$NP$ と $k$ の交点を $Q$ とする。

$AB$ 上を $P$ がうごく時, $Q$ は $C, D$ を両端とする大円の劣弧上を動く。この長さは, 半径 1 の四分円の弧長にひとしく,
\[
\frac{1}{4} \cdot 2\pi = \frac{1}{2}\pi \quad \cdots \text{①}
\]

次に $AB$ 上を $P$ がうごく時, 平面 $ABN$ で $k$ を切ると下図で $\text{+++}$ 上を $Q$ がくまなくうごく。この円は, 1 辺の長さ $2\sqrt{2}$ の正三角形の外接円で, その半径 $\frac{2\sqrt{6}}{3}$ だから
\[
\stackrel{\frown}{CD} = \frac{1}{3} \cdot \frac{2\sqrt{6}}{3} \pi = \frac{2\sqrt{6}}{9}\pi \quad \cdots \text{②}
\]

\begin{center}
\begin{tikzpicture}[scale=1.3]
  % Diagram 1: Sphere sketch
  \begin{scope}[shift={(0,0)}]
    \draw (0,0) circle (1.2);
    \draw (-1.2,0) arc (180:360:1.2 and 0.4);
    \draw[dashed] (-1.2,0) arc (180:0:1.2 and 0.4);
    
    \coordinate (N) at (0,1.2);
    \coordinate (A) at (-1.8,-1.5);
    \coordinate (B) at (1.8,-1.5);
    
    \draw (N) -- (A);
    \draw (N) -- (B);
    \draw (A) -- (B);
    
    \node[above] at (N) {$N$};
    \node[below left] at (A) {$A$};
    \node[below right] at (B) {$B$};
  \end{scope}

  % Diagram 2: Triangle ABN section
  \begin{scope}[shift={(5,0)}]
    \coordinate (N) at (0, 2.45);
    \coordinate (A) at (-2, 0);
    \coordinate (B) at (2, 0);

    \draw[thick] (N) -- (A) -- (B) -- cycle;

    \coordinate (C) at ($(N)!0.5!(A)$);
    \coordinate (D) at ($(N)!0.5!(B)$);

    % Outer circle / arc CD
    \draw[thick] (C) arc (210:330:2.3);

    \node[above] at (N) {$N$};
    \node[below left] at (A) {$A$};
    \node[below right] at (B) {$B$};
    \node[left] at (C) {$C$};
    \node[right] at (D) {$D$};

    \node[left] at ($(N)!0.5!(C)$) {$\sqrt{2}$};
    \node[left] at ($(C)!0.5!(A)$) {$\sqrt{2}$};
    \node[right] at ($(N)!0.5!(D)$) {$\sqrt{2}$};
    \node[right] at ($(D)!0.5!(B)$) {$2\sqrt{2}$};
    \node[below] at ($(A)!0.5!(B)$) {$2\sqrt{2}$};

    \node[below] at (0,1.0) {$\frac{2}{3}\pi$};
  \end{scope}
\end{tikzpicture}
\end{center}

以上の ①, ② から, もとめる長さは, (これらが $C, D$ 以外で交わらないことから)
\[
\frac{1}{2}\pi + \frac{2\sqrt{6}}{9}\pi = \left(\frac{1}{2} + \frac{2\sqrt{6}}{9}\right)\pi
\]

\end{document}
